\begin{ouverture}[courbe]{9}{Programmes de calcul et dépendance}{Partie 2 · Algèbre et modèles}
{Programmes · Formules · Tableaux · Graphiques}
\objectif{Appliquer un programme de calcul à un nombre, puis à une lettre.}
\objectif{Remonter un programme pour retrouver le nombre de départ.}
\objectif{Traduire la dépendance de deux grandeurs par une formule, un tableau et un graphique.}
\prerequis{Repérage dans le plan, calcul littéral et équations.}
\end{ouverture}
\pagestyle{otMain}\essentiel\fontsize{11}{15}\selectfont
\phantomsection\label{cours-9-1}
\essSection{1}{Des programmes de calcul aux formules}
\begin{definition}Un programme de calcul enchaîne des opérations à partir d’un nombre de départ. Appliqué à une lettre, il produit une expression littérale qui donne le résultat pour tous les nombres de départ à la fois.\end{definition}
\begin{exemple}Programme : choisir un nombre ; soustraire 4 ; multiplier le résultat par 3. Avec 10 : $10-4=6$, puis $6\times3=18$. Avec $x$ : $3(x-4)=3x-12$. On retrouve 18 pour $x=10$ : $3\times10-12=18$.\end{exemple}
\begin{savoirfaire}[Remonter un programme]Pour retrouver le nombre de départ qui donne 9, on effectue les opérations inverses dans l’ordre inverse : $9:3=3$, puis $3+4=7$. Contrôle : $(7-4)\times3=9$. On peut aussi résoudre l’équation $3(x-4)=9$.\end{savoirfaire}
\begin{definition}[Une notation]Quand une grandeur $P$ dépend d’une grandeur $t$, on note $P(t)$ sa valeur pour une valeur donnée de $t$. Si $P(t)=6+2t$ est le prix en euros d’une location de $t$ heures, $P(3)=12$ se lit « $P$ de 3 égale 12 » : trois heures coûtent 12 €.\end{definition}
\begin{attention}$P(3)$ n’est pas un produit : c’est la valeur de $P$ pour $t=3$. Le programme et la formule décrivent la même dépendance ; le tableau et le graphique en montrent quelques valeurs.\end{attention}
\clearpage
\phantomsection\label{cours-9-2}
\essSection{2}{Tableaux et graphiques}
\begin{exemple}Pour $P(t)=6+2t$ : $P(0)=6$, $P(1)=8$, $P(2)=10$, $P(3)=12$. Les points de coordonnées $(t\,;P(t))$ sont placés dans un repère, la durée en abscisse et le prix en ordonnée.\end{exemple}
\label{fig-cours-9-2-cours-location}\figureGeo{cours-location}

\begin{savoirfaire}[Lire dans les deux sens]Partir d’une valeur de l’axe horizontal, rejoindre le tracé verticalement, puis lire l’axe vertical. Dans l’autre sens, suivre l’horizontale d’une valeur verticale jusqu’à toutes ses intersections avec le tracé.\end{savoirfaire}
\begin{propriete}Une situation de proportionnalité se représente par des points alignés avec l’origine. Ici $P(0)=6$ : le prix n’est pas proportionnel à la durée. L’aire $A(c)=c^2$ d’un carré de côté $c$ se représente par une courbe : ce n’est pas une proportionnalité.\end{propriete}
\begin{attention}Lire d’abord les graduations et les unités. Une lecture graphique est souvent approchée et ne vaut que dans le domaine représenté.\end{attention}

\clearpage\rubrique{Découverte}
\begin{blocLivre}{ex-9-1}\begin{exercice}[1]{Appliquer un programme}
Programme : choisir un nombre ; le multiplier par 3 ; soustraire 5 au résultat. Appliquer ce programme à 4, à $-2$ et à 1,5.
\end{exercice}\end{blocLivre}
\begin{blocLivre}{ex-9-2}\begin{exercice}[1]{Du programme à l’expression}
Programme : choisir un nombre ; lui ajouter 6 ; multiplier le résultat par $-2$. Écrire le résultat obtenu à partir de $x$, puis développer.
\end{exercice}\end{blocLivre}
\begin{blocLivre}{ex-9-3}\begin{exercice}[1]{Remonter un programme}
Programme : choisir un nombre ; soustraire 4 ; multiplier le résultat par 5. On obtient 35. Retrouver le nombre de départ de deux façons.
\end{exercice}\end{blocLivre}
\begin{blocLivre}{ex-9-4}\begin{exercice}[1]{Une formule de tarif}
Une location coûte 5 euros fixes et 3 euros par heure. Exprimer le prix $P(t)$ pour $t$ heures, puis calculer $P(0)$, $P(2)$ et $P(4)$.
\end{exercice}\end{blocLivre}
\begin{blocLivre}{ex-9-5}\begin{exercice}[1]{Lire un couple}
Le graphique donne la température $T$ d’un four, en $^\circ$C, en fonction du temps en minutes. Lire $T(2)$.
\label{fig-ex-9-5}\figureGeo{four}
\end{exercice}\end{blocLivre}

\rubrique{Entraînement}
\begin{blocLivre}{ex-9-6}\begin{exercice}[2]{Lire une durée}
Sur le même graphique, à quel instant la température atteint-elle $100\,^\circ$C pour la première fois ?
\label{fig-ex-9-6}\figureGeo{four}
\end{exercice}\end{blocLivre}
\begin{blocLivre}{ex-9-7}\begin{exercice}[2]{Une lecture inverse double}
Sur le graphique du four, donner tous les instants où la température vaut $60\,^\circ$C.
\label{fig-ex-9-7}\figureGeo{four}
\end{exercice}\end{blocLivre}
\begin{blocLivre}{ex-9-8}\begin{exercice}[2]{Reconnaître une relation}
Le graphique représente-t-il une situation de proportionnalité ? Justifier.
\label{fig-ex-9-8}\figureGeo{graph-non-proportion}
\end{exercice}\end{blocLivre}
\begin{blocLivre}{ex-9-9}\begin{exercice}[2]{Construire un graphique}
On vend une corde à 1,5 euro le mètre, pour des longueurs de 0 à 6 m. Compléter un tableau pour 0, 2, 4 et 6 m, puis tracer la représentation du prix en fonction de la longueur.
\label{fig-ex-9-9}\figureGeo{graph-corde}
\end{exercice}\end{blocLivre}
\begin{blocLivre}{ex-9-10}\begin{exercice}[2]{Aire d’un carré}
Exprimer l’aire $A(c)$ d’un carré de côté $c$ cm. Calculer $A(1)$, $A(2)$, $A(3)$ et $A(4)$. L’aire est-elle proportionnelle au côté ?
\label{fig-ex-9-10}\figureGeo{graph-carre}
\end{exercice}\end{blocLivre}

\rubrique{Problèmes}
\begin{blocLivre}{ex-9-11}\begin{exercice}[3]{Deux programmes égaux}
Programme A : choisir un nombre ; ajouter 3 ; multiplier par 4 ; soustraire 12. Programme B : choisir un nombre ; le multiplier par 4. Montrer qu’ils donnent toujours le même résultat. Quel nombre de départ donne 26 ?
\end{exercice}\end{blocLivre}
\begin{blocLivre}{ex-9-12}\begin{exercice}[3]{Un tableur}
Dans un tableur, A2 contient une durée en heures. Quelle formule saisir en B2 pour calculer le prix $P(t)=5+3t$, puis recopier vers le bas ? Quels prix s’affichent en B3 et B4 ?
\label{fig-ex-9-12}\figureGeo{tableur-prix}
\end{exercice}\end{blocLivre}
\begin{blocLivre}{ex-9-13}\begin{exercice}[3]{Retirer des carrés}
Dans un rectangle de 18 cm sur 10 cm, on retire aux quatre coins des carrés de côté $x$ cm, avec $0<x<5$. Exprimer l’aire restante $A(x)$, puis calculer $A(2)$.
\label{fig-ex-9-13}\figureGeo{coins-retires}
\end{exercice}\end{blocLivre}
\begin{blocLivre}{ex-9-14}\begin{exercice}[3]{Deux courbes de prix}
Les graphiques correspondent aux prix $A(t)=4+2t$ et $B(t)=4t$, pour $0\le t\le4$ heures. Lire l’heure d’égalité, puis la vérifier par une équation.
\label{fig-ex-9-14}\figureGeo{graph-tarifs}
\end{exercice}\end{blocLivre}
\begin{blocLivre}{ex-9-15}\begin{exercice}[3]{Une réserve qui diminue}
Une réserve contient 120 L ; on retire 8 L par minute à débit constant. Donner une formule du volume $V(t)$ après $t$ minutes et le domaine de validité de ce modèle.
\label{fig-ex-9-15}\figureGeo{graph-reserve}
\end{exercice}\end{blocLivre}

\clearpage
